% GENERATED FILE. Edit canonical lessons, then run npm run generate:books.
% canonical-source-hash: fnv1a64:8317b9c4f600a325
% canonical-lessons: GE-C291-notruf, GE-C291-notarzt, GE-C291-hubschrauber, GE-C291-sanitaeter, GE-C291-notaufnahme

\chapter{Der Notruf}
\label{ch:ge-der-notruf}
\hlchaptermodality{291}

\begin{chapteropening}
\textbf{By the end of this chapter:} \emph{I can describe calling for help: the emergency number, the emergency doctor, a helicopter, a paramedic and the emergency room.}

Complete the last lesson of chapter 291: I can describe calling for help: the emergency number, the emergency doctor, a helicopter, a paramedic and the emergency room.
\end{chapteropening}

\section[der Notruf]{der Notruf — the emergency call; the emergency number}

\label{lesson:GE-C291-notruf}



\begin{quote}
Before the new one: say the German for to choke, then the German for to drown.
\end{quote}

\subsection*{You'll want to know: der Notruf}
\textbf{der Notruf} — \textquotedblleft{}the emergency call; the emergency number\textquotedblright{}.

\begin{quote}
Sie wählte sofort den Notruf.
\end{quote}

\begin{quote}
\emph{She dialled the emergency number at once.}
\end{quote}

Plural: \emph{die Notrufe}. The number 112 reaches the fire brigade and the ambulance.

\begin{cousinweb}
\emph{Notruf} joins \emph{Not} \textquotedblleft{}need\textquotedblright{}, a cousin of English \emph{need}, and \emph{Ruf} \textquotedblleft{}call\textquotedblright{}, from \emph{rufen}, Proto-Germanic \emph{hrōpaną} (EtymDB).
\end{cousinweb}

\begin{grammarlens}[title={What you've built}]
One more word for telling, explaining and describing at B1.
\end{grammarlens}

\subsection*{Your turn}
\begin{itemize}
\raggedright
  \item \emph{Say it:} \emph{der Notruf}
  \item \emph{Say it:} \emph{der Notruf}, once more
  \item \emph{Say it:} a sentence of your own with \emph{der Notruf}
  \item \emph{Recall:} say the German for to choke, then the German for to drown, then say \emph{der Notruf} again
\end{itemize}

\begin{tcolorbox}[breakable,colback=teal!4,colframe=teal!35!black,title={Before you move on}]
What does der Notruf mean? (\textquotedblleft{}The emergency call; the emergency number\textquotedblright{}.) Say it once more.
\end{tcolorbox}
\section[der Notarzt]{der Notarzt — the emergency doctor}

\label{lesson:GE-C291-notarzt}



\begin{quote}
Before the new one: say the German for to drown, then the German for the emergency call.
\end{quote}

\subsection*{You'll want to know: der Notarzt}
\textbf{der Notarzt} — \textquotedblleft{}the emergency doctor\textquotedblright{}.

\begin{quote}
Der Notarzt kam nach zehn Minuten.
\end{quote}

\begin{quote}
\emph{The emergency doctor came after ten minutes.}
\end{quote}

Plural: \emph{die Notärzte}. Feminine \emph{die Notärztin}.

\begin{cousinweb}
\emph{Notarzt} joins \emph{Not} \textquotedblleft{}need\textquotedblright{} and \emph{Arzt} \textquotedblleft{}doctor\textquotedblright{}, from Old High German \emph{arzat}, which came from Latin \emph{archiāter} \textquotedblleft{}physician\textquotedblright{} (EtymDB).
\end{cousinweb}

\begin{grammarlens}[title={What you've built}]
One more word for telling, explaining and describing at B1.
\end{grammarlens}

\subsection*{Your turn}
\begin{itemize}
\raggedright
  \item \emph{Say it:} \emph{der Notarzt}
  \item \emph{Say it:} \emph{der Notarzt}, once more
  \item \emph{Say it:} a sentence of your own with \emph{der Notarzt}
  \item \emph{Recall:} say the German for to drown, then the German for the emergency call, then say \emph{der Notarzt} again
\end{itemize}

\begin{tcolorbox}[breakable,colback=teal!4,colframe=teal!35!black,title={Before you move on}]
What does der Notarzt mean? (\textquotedblleft{}The emergency doctor\textquotedblright{}.) Say it once more.
\end{tcolorbox}
\section[der Hubschrauber]{der Hubschrauber — the helicopter}

\label{lesson:GE-C291-hubschrauber}



\begin{quote}
Before the new one: say the German for the emergency call, then the German for the emergency doctor.
\end{quote}

\subsection*{You'll want to know: der Hubschrauber}
\textbf{der Hubschrauber} — \textquotedblleft{}the helicopter\textquotedblright{}.

\begin{quote}
Ein Hubschrauber brachte die Frau ins Krankenhaus.
\end{quote}

\begin{quote}
\emph{A helicopter took the woman to hospital.}
\end{quote}

Plural: \emph{die Hubschrauber}. Built from \emph{Hub} 'lift' and \emph{Schraube} 'screw'.

\begin{cousinweb}
\emph{Hubschrauber} is a modern compound of \emph{Hub} \textquotedblleft{}a lift\textquotedblright{}, from \emph{heben} \textquotedblleft{}to lift\textquotedblright{}, and \emph{Schrauber}, from \emph{Schraube} \textquotedblleft{}screw\textquotedblright{}; EtymDB lists no parent: a screw that lifts.
\end{cousinweb}

\begin{grammarlens}[title={What you've built}]
One more word for telling, explaining and describing at B1.
\end{grammarlens}

\subsection*{Your turn}
\begin{itemize}
\raggedright
  \item \emph{Say it:} \emph{der Hubschrauber}
  \item \emph{Say it:} \emph{der Hubschrauber}, once more
  \item \emph{Say it:} a sentence of your own with \emph{der Hubschrauber}
  \item \emph{Recall:} say the German for the emergency call, then the German for the emergency doctor, then say \emph{der Hubschrauber} again
\end{itemize}

\begin{tcolorbox}[breakable,colback=teal!4,colframe=teal!35!black,title={Before you move on}]
What does der Hubschrauber mean? (\textquotedblleft{}The helicopter\textquotedblright{}.) Say it once more.
\end{tcolorbox}
\section[der Sanitäter]{der Sanitäter — the paramedic}

\label{lesson:GE-C291-sanitaeter}



\begin{quote}
Before the new one: say the German for the emergency doctor, then the German for the helicopter.
\end{quote}

\subsection*{You'll want to know: der Sanitäter}
\textbf{der Sanitäter} — \textquotedblleft{}the paramedic\textquotedblright{}.

\begin{quote}
Die Sanitäter halfen dem Verletzten.
\end{quote}

\begin{quote}
\emph{The paramedics helped the injured man.}
\end{quote}

Plural: \emph{die Sanitäter}. Colloquially \emph{die Sanis}.

\begin{cousinweb}
\emph{Sanitäter} is built on \emph{Sanität} \textquotedblleft{}health service\textquotedblright{}, a learned borrowing from Latin \emph{sānitās} \textquotedblleft{}health\textquotedblright{}, from \emph{sānus} \textquotedblleft{}healthy\textquotedblright{} (Lewis \& Short).
\end{cousinweb}

\begin{grammarlens}[title={What you've built}]
One more word for telling, explaining and describing at B1.
\end{grammarlens}

\subsection*{Your turn}
\begin{itemize}
\raggedright
  \item \emph{Say it:} \emph{der Sanitäter}
  \item \emph{Say it:} \emph{der Sanitäter}, once more
  \item \emph{Say it:} a sentence of your own with \emph{der Sanitäter}
  \item \emph{Recall:} say the German for the emergency doctor, then the German for the helicopter, then say \emph{der Sanitäter} again
\end{itemize}

\begin{tcolorbox}[breakable,colback=teal!4,colframe=teal!35!black,title={Before you move on}]
What does der Sanitäter mean? (\textquotedblleft{}The paramedic\textquotedblright{}.) Say it once more.
\end{tcolorbox}
\section[die Notaufnahme]{die Notaufnahme — the emergency room, A and E}

\label{lesson:GE-C291-notaufnahme}



\begin{quote}
Before the new one: say the German for the helicopter, then the German for the paramedic.
\end{quote}

\subsection*{You'll want to know: die Notaufnahme}
\textbf{die Notaufnahme} — \textquotedblleft{}the emergency room, A and E\textquotedblright{}.

\begin{quote}
In der Notaufnahme musste er zwei Stunden warten.
\end{quote}

\begin{quote}
\emph{In the emergency room he had to wait two hours.}
\end{quote}

Plural: \emph{die Notaufnahmen}. \emph{Die Aufnahme} — admission.

\begin{cousinweb}
\emph{Notaufnahme} joins \emph{Not} \textquotedblleft{}need\textquotedblright{} and \emph{Aufnahme} \textquotedblleft{}admission\textquotedblright{}, from \emph{aufnehmen} \textquotedblleft{}to take in\textquotedblright{}, built on \emph{nehmen}, Proto-Germanic \emph{nemaną} (EtymDB).
\end{cousinweb}

\begin{grammarlens}[title={What you've built}]
One more word for telling, explaining and describing at B1.
\end{grammarlens}

\subsection*{Your turn}
\begin{itemize}
\raggedright
  \item \emph{Say it:} \emph{die Notaufnahme}
  \item \emph{Say it:} \emph{die Notaufnahme}, once more
  \item \emph{Say it:} a sentence of your own with \emph{die Notaufnahme}
  \item \emph{Recall:} say the German for the helicopter, then the German for the paramedic, then say \emph{die Notaufnahme} again
\end{itemize}

\begin{tcolorbox}[breakable,colback=teal!4,colframe=teal!35!black,title={Before you move on}]
What does die Notaufnahme mean? (\textquotedblleft{}The emergency room, A and E\textquotedblright{}.) Say it once more.
\end{tcolorbox}
